\documentclass[10pt,a4paper]{article}

\usepackage[T1]{fontenc}
\usepackage[utf8]{inputenc}
\usepackage[ngerman]{babel}
\usepackage[a4paper,margin=14mm]{geometry}
\usepackage{lmodern}
\usepackage[table]{xcolor}
\usepackage{tabularx}
\usepackage{array}
\usepackage{microtype}
\usepackage[most]{tcolorbox}
\usepackage{tikz}
\usetikzlibrary{arrows.meta,shapes.geometric}
\usepackage{textcomp}

% Gerade Anführungszeichen; " ist bei ngerman ein Kurzbefehl.
\newcommand{\qd}[1]{\textquotedbl#1\textquotedbl}

\pagestyle{empty}
\setlength{\parindent}{0pt}
\setlength{\parskip}{0pt}
\renewcommand{\familydefault}{\sfdefault}

\definecolor{Navy}{HTML}{17324D}
\definecolor{Blue}{HTML}{3B6EA5}
\definecolor{Sky}{HTML}{EAF3FA}
\definecolor{Mint}{HTML}{E8F5EF}
\definecolor{Gold}{HTML}{F3B33D}
\definecolor{Ink}{HTML}{1E2935}
\definecolor{Muted}{HTML}{617080}
\definecolor{Line}{HTML}{D8E1E8}
\definecolor{CodeBg}{HTML}{F5F7F9}
\definecolor{Cell}{HTML}{FFF4DB}

\color{Ink}

\newcommand{\sectiontitle}[1]{%
  \vspace{2.3mm}
  {\color{Blue}\bfseries\large #1}\par
  \vspace{0.8mm}
  {\color{Blue}\rule{\linewidth}{0.7pt}}\par
  \vspace{1.2mm}
}

\tikzset{
  func/.style={draw=Ink, line width=0.7pt, ellipse, fill=white,
               minimum width=22mm, minimum height=8mm, inner sep=1pt,
               font=\small},
  io/.style={font=\small, inner sep=1pt},
  hint/.style={font=\scriptsize\bfseries, text=Blue, inner sep=1pt},
  flow/.style={draw=Ink, line width=0.7pt,
               -{Stealth[length=2.2mm,width=1.8mm]}},
}

\begin{document}

\colorbox{Navy}{%
  \parbox{\dimexpr\linewidth-2\fboxsep\relax}{%
    \vspace{3mm}\color{white}
    \begin{tabularx}{\linewidth}{@{}Xr@{}}
      {\small\bfseries INFORMATIK · KLASSE 9} &
      \colorbox{Gold}{\color{Navy}\small\bfseries\strut LÖSUNGEN}
    \end{tabularx}
    \vspace{2.5mm}

    {\fontsize{22}{25}\selectfont\bfseries Sicherungsblatt: Aufgabe 5b}\par
    \vspace{1mm}
    {\large WENN-Funktion als Datenflussdiagramm}
    \vspace{3mm}
  }%
}

\vspace{3mm}
\colorbox{Sky}{%
  \parbox{\dimexpr\linewidth-2\fboxsep\relax}{%
    \vspace{1.8mm}
    \textcolor{Blue}{\bfseries Ziel:}
    Du stellst \textbf{Aussagefunktion} und \textbf{WENN-Funktion} als
    Datenflussdiagramm dar und findest einen Fehler darin.
    \vspace{1.8mm}
  }%
}

\sectiontitle{1 · Aussagefunktion}

\begin{tikzpicture}[x=1mm,y=1mm,
    desc/.style={font=\small, text=Ink, anchor=north west, inner sep=0pt,
                 text width=62mm, align=left}]

  \node[io] (note) at (16,30) {Note};
  \node[io] (zwei) at (44,30) {2};
  \node[func] (kleiner) at (30,17) {kleiner};
  \draw[flow] (note.south) -- (kleiner.north west);
  \draw[flow] (zwei.south) -- (kleiner.north east);
  \node[io] (wert) at (30,3) {Wahrheitswert};
  \draw[flow] (kleiner.south) -- (wert.north);
  \node[hint] at (30,-1) {WAHR oder FALSCH};

  \node[desc] at (62,29) {Eine \textbf{Aussagefunktion} vergleicht Werte. Ihr
    Ergebnis ist ein \textbf{Wahrheitswert}: WAHR oder FALSCH.\\[1.5mm]
    Weitere: größer (\texttt{>}), größer gleich (\texttt{>=}), gleich
    (\texttt{=}), ISTGERADE (nur ein Eingang).};

  \node[anchor=north west, inner sep=0pt] at (132,31) {%
    \renewcommand{\arraystretch}{1.1}%
    \begin{tabular}{@{}>{\centering\arraybackslash\small}p{9mm}>{\centering\arraybackslash\small}p{12mm}>{\centering\arraybackslash\small}p{16mm}@{}}
      \rowcolor{Sky}\bfseries Note & \bfseries Vergleich & \bfseries Ergebnis \\
      1 & $1 < 2$ & WAHR \\
      \rowcolor{CodeBg}2 & $2 < 2$ & FALSCH \\
      3 & $3 < 2$ & FALSCH \\
    \end{tabular}};
  \node[anchor=north west, font=\footnotesize, text=Muted, text width=42mm,
        align=left, inner sep=0pt] at (132,11)
    {2 ist nicht kleiner als 2.};
\end{tikzpicture}

\sectiontitle{2 · WENN-Funktion: Notengeld aus Aufgabe 4}

\begin{tikzpicture}[x=1mm,y=1mm,
    desc/.style={font=\small, text=Ink, anchor=north west, inner sep=0pt,
                 text width=64mm, align=left}]

  \node[io] (note) at (16,48) {Note};
  \node[io] (zwei) at (44,48) {2};
  \node[func] (kleiner) at (30,38) {kleiner};
  \draw[flow] (note.south) -- (kleiner.north west);
  \draw[flow] (zwei.south) -- (kleiner.north east);

  \node[hint] at (72,33) {2 · Ja-Fall};
  \node[io] (ja) at (72,28.5) {\qd{5\texteuro}};
  \node[hint] at (100,33) {3 · Nein-Fall};
  \node[io] (nein) at (100,28.5) {\qd{0\texteuro}};

  \node[func] (wenn) at (62,15) {WENN};
  \draw[flow] (kleiner.south) -- (wenn.north west);
  \node[hint, anchor=east] at (38.5,26) {1 · Wahrheitswert};
  \draw[flow] (ja.south) -- (wenn.north);
  \draw[flow] (nein.south) -- (wenn.north east);

  \node[io] (geld) at (62,2) {Notengeld};
  \draw[flow] (wenn.south) -- (geld.north);

  \node[desc] at (118,50) {Die \textbf{WENN-Funktion} hat drei Eingänge:\\[1mm]
    \textbf{1.}~den Wahrheitswert einer Aussagefunktion,\\
    \textbf{2.}~den Wert für den Ja-Fall,\\
    \textbf{3.}~den Wert für den Nein-Fall.\\[1.5mm]
    Note 1 \textrightarrow{} WAHR \textrightarrow{} \qd{5\texteuro}\\
    Note 3 \textrightarrow{} FALSCH \textrightarrow{} \qd{0\texteuro}\\[1.5mm]
    \textcolor{Blue}{\bfseries Termdarstellung:}\\[0.5mm]
    \colorbox{CodeBg}{\ttfamily\small =WENN(Note<2; \qd{5\texteuro}; \qd{0\texteuro})}};
\end{tikzpicture}

\sectiontitle{3 · Fehler finden: Gewinnspiel}

{\small Gewinn 10\,\texteuro{} nur bei Losnummer \textbf{größer als 70} --
im Diagramm stand \textbf{größer gleich}.\par}

\vspace{1.5mm}
\begin{tikzpicture}[x=1mm,y=1mm,
    desc/.style={font=\small, text=Ink, anchor=north west, inner sep=0pt,
                 text width=64mm, align=left}]

  \node[io] (los) at (16,40) {Losnummer};
  \node[io] (siebzig) at (44,40) {70};
  \node[func, draw=Blue, line width=1.2pt] (groesser) at (30,30) {größer};
  \draw[flow] (los.south) -- (groesser.north west);
  \draw[flow] (siebzig.south) -- (groesser.north east);
  \node[hint, anchor=west] at (43,26) {korrigiert: \texttt{>} statt \texttt{>=}};

  \node[io] (ja) at (80,25) {\qd{10\texteuro}};
  \node[io] (nein) at (104,25) {\qd{0\texteuro}};
  \node[func] (wenn) at (70,11) {WENN};
  \draw[flow] (groesser.south) -- (wenn.north west);
  \draw[flow] (ja.south) -- (wenn.north);
  \draw[flow] (nein.south) -- (wenn.north east);
  \node[io] (gewinn) at (70,-1) {Losgewinn};
  \draw[flow] (wenn.south) -- (gewinn.north);

  \node[anchor=north west, inner sep=0pt] at (118,42) {%
    \renewcommand{\arraystretch}{1.1}%
    \begin{tabular}{@{}>{\centering\arraybackslash\small}p{16mm}>{\centering\arraybackslash\small}p{20mm}>{\centering\arraybackslash\small}p{20mm}@{}}
      \rowcolor{Sky}\bfseries Losnummer & \bfseries mit \texttt{>=} & \bfseries mit \texttt{>} \\
      \cellcolor{Cell}70 & \cellcolor{Cell}10\,\texteuro & \cellcolor{Cell}0\,\texteuro \\
      71 & 10\,\texteuro & 10\,\texteuro \\
      \rowcolor{CodeBg}90 & 10\,\texteuro & 10\,\texteuro \\
    \end{tabular}};
  \node[desc] at (118,20.5) {Der Fehler wirkt sich nur bei der Losnummer 70
    aus.\\[1mm]
    \colorbox{CodeBg}{\ttfamily\small =WENN(Losnummer>70; \qd{10\texteuro}; \qd{0\texteuro})}};
\end{tikzpicture}

\sectiontitle{4 · Merke}

\begin{tcolorbox}[colback=Mint,colframe=Blue!40,arc=3mm,boxrule=0.7pt,
  left=3mm,right=3mm,top=2mm,bottom=2mm]
  Eine \textbf{Aussagefunktion} vergleicht Werte und gibt einen
  \textbf{Wahrheitswert} aus: WAHR oder FALSCH.\par
  \vspace{1.2mm}
  Die \textbf{WENN-Funktion} hat drei Eingänge: Wahrheitswert, Ja-Fall,
  Nein-Fall. Sie gibt genau einen der beiden Werte aus.
\end{tcolorbox}

\vfill
{\color{Line}\rule{\linewidth}{0.5pt}}\par
\vspace{1.5mm}
\begin{tabularx}{\linewidth}{@{}Xr@{}}
  {\small\color{Muted}Klasse 9 · Tabellenkalkulation} &
  {\small\color{Muted}Sicherungsblatt · Aufgabe 5b · Version 2}
\end{tabularx}

\end{document}
